% ==================================================================
% §15.5.2 REBUILT — Phase 1: Retained five-probe effective analysis
%
% This block is Part 2 of the §15.5 rebuild under the parallel-creation
% strategy (PROJECT_STRUCTURE.md §3.6). It is inserted IMMEDIATELY AFTER
% the §15.5.1 block (label: subsec:eps_def_rebuilt) and BEFORE the
% closing delimiter of the rebuild parent block.
%
% LOCKED 5-column retained-probe table schema (PROJECT_STRUCTURE.md §3.8).
% Numerics from scripts/epsilon_convergence/result_ch{1..4,8}.json.
% GPT v2 round 5 corrections applied (2026-04-18).
% ==================================================================

\subsubsection*{\normalfont\textbf{15.5.2 \quad Retained five-probe effective analysis}}
\label{subsec:eps_retained_rebuilt}

With the effective meaning of $\varepsilon$ fixed in
\S\ref{subsec:eps_def_rebuilt}, we now specify which cosmological probes
are retained for the present uniform-$\varepsilon$ analysis, what each
of them actually constrains, and how their resulting effective intervals
are to be interpreted.

\paragraph{Retention criterion.}
A probe is retained in \S15.5 only if it constrains the present
effective $\varepsilon$-sector rather than a distinct physical sector,
and if its extraction pipeline is sufficiently explicit that an
effective interval can be stated without importing an uncontrolled
external likelihood. Five probes satisfy this criterion under the
uniform-$\varepsilon$ ansatz and the $\Lambda$CDM-background proxy
$(\Omega_m = 0.315,\ \Omega_\Lambda = 0.685)$; all other probes
considered in the present analysis are discussed in \S15.5.4.

\begin{table}[htbp]
\centering
\renewcommand{\arraystretch}{1.25}
\small
\begin{tabular}{@{}p{2.6cm}p{3.1cm}p{3.4cm}p{3.0cm}p{1.6cm}@{}}
\toprule
\textbf{Probe} & \textbf{Nature of constraint} & \textbf{What is actually fitted} & \textbf{Kernel / epoch sensitivity} & \textbf{Role} \\
\midrule
Hubble $+\,r_s$ &
two-sided effective extraction &
$\varepsilon$ from the Hubble-channel distance-ratio proxy under the present simplified $r_s$ treatment &
broad integrated kernel, CMB-era proxy scale to $z=0$ &
Primary \\
\addlinespace
JWST early-galaxy excess &
two-sided effective extraction (broad) &
$\varepsilon$ from the Press--Schechter proxy for the high-$z$ abundance enhancement at $z\approx 8$--$12$ &
high-$z$ structure-growth window &
Primary \\
\addlinespace
Cosmic chronometers &
one-sided consistency bound under $\varepsilon\geq 0$ &
$\chi^2$ fit of $H(z)$ residuals in the $\Lambda$CDM-background proxy &
direct late-time expansion, $0.07\leq z\leq 1.97$ &
Consistency \\
\addlinespace
$f\sigma_8$ from RSD &
one-sided consistency bound under $\varepsilon\geq 0$ &
$\varepsilon$ in the linear-growth ODE likelihood &
late-time growth, $z\lesssim 2$ &
Consistency \\
\addlinespace
ISW amplitude &
one-sided consistency bound under $\varepsilon\geq 0$ &
$\varepsilon$ through the linearised ISW proxy amplitude &
late-time potential evolution, $z\lesssim 2$ &
Consistency \\
\bottomrule
\end{tabular}
\caption{Retained five-probe set for the present effective
uniform-$\varepsilon$ analysis layer. ``Primary'' probes provide two-sided
effective extraction intervals; ``Consistency'' probes provide
one-sided upper bounds under the physical prior
$\varepsilon\geq 0$ introduced in \S\ref{subsec:eps_def_rebuilt}.}
\label{tab:eps_retained_probes_rebuilt}
\end{table}

\paragraph{Two-sided vs upper-bound distinction.}
Within the retained five-probe set, Hubble$+r_s$ and the JWST
early-galaxy excess provide two-sided effective extraction intervals,
whereas cosmic chronometers, $f\sigma_8$ from RSD, and the ISW
amplitude are retained as late-time consistency channels whose raw
likelihoods are compatible with $\varepsilon = 0$ and are therefore
reported here as one-sided upper bounds under the physical prior
$\varepsilon\geq 0$ introduced in \S\ref{subsec:eps_def_rebuilt}.
In the one-sided cases quoted below, the reported $1\sigma$ and
$2\sigma$ bounds are the upper edges of the corresponding raw
intervals after imposing the physical prior $\varepsilon \geq 0$;
they are not separate positive-domain refits.

\paragraph{Probe 1 --- Hubble $+\,r_s$ (two-sided effective extraction).}
The SH0ES distance-ladder determination
$H_0 = 73.04\pm1.04~\mathrm{km\,s^{-1}\,Mpc^{-1}}$~\cite{Riess2022}
and the Planck 2018 base-$\Lambda$CDM value
$H_0 = 67.4\pm0.5~\mathrm{km\,s^{-1}\,Mpc^{-1}}$~\cite{Planck2018}
define an observed tension of $\sim 5\sigma$.
The effective distance-ratio proxy adopted here for the Hubble-tension
channel, evaluated within the uniform-$\varepsilon$ ansatz on the
$\Lambda$CDM-background proxy and with the present simplified treatment
of $r_s$, yields a two-sided effective interval centred at
\begin{equation}
  \varepsilon_{\rm H} \;=\; 0.032,
  \qquad
  \text{1\,\ensuremath{\sigma}:}\ [0.027,\,0.038],
  \qquad
  \text{2\,\ensuremath{\sigma}:}\ [0.021,\,0.043].
  \label{eq:eps_hubble_retained}
\end{equation}
This should be read as an effective H1-type extraction under the
present $r_s$ treatment, not as a full Boltzmann/CMB-likelihood
implementation of the H2 programme.

\paragraph{Probe 2 --- JWST early-galaxy excess (two-sided effective extraction; binding below).}
Reported high-$z$ stellar-mass abundances at $z\approx 8$--$12$ suggest
a representative enhancement target of order $R \sim 10$, with a broader
observationally discussed range of roughly $3$--$100$, relative to
baseline $\Lambda$CDM expectations for the relevant mass scales. A
Press--Schechter proxy tying $G_{\rm eff}(z)$ to the halo mass function
through the linear growth factor, on the same $\Lambda$CDM-background
proxy and without full growth-ODE integration or halo-to-stellar-mass
mapping, is translated here into a two-sided effective interval centred at
\begin{equation}
  \varepsilon_{\rm JWST} \;=\; 0.042,
  \qquad
  \text{1\,\ensuremath{\sigma}:}\ [0.029,\,0.071],
  \qquad
  \text{2\,\ensuremath{\sigma}:}\ [0.013,\,0.179].
  \label{eq:eps_jwst_retained}
\end{equation}
This interval should be read as a proxy-level effective extraction
within the present Press--Schechter treatment, not as the result of a
full structure-formation pipeline with transfer-function, baryonic, and
halo-to-stellar-mass systematics propagated end-to-end. In the present
five-probe comparison, the lower edge $\varepsilon = 0.029$ will
become the lower edge of the joint effective band discussed in
\S15.5.3.

\paragraph{Probe 3 --- Cosmic chronometers (one-sided consistency bound under $\varepsilon\geq 0$).}
Moresco-compilation $H(z)$ points at $z=0.07$--$1.97$ are fit against
the effective expansion
$H^{\rm ECT}(z) = H_0\sqrt{\Omega_m(1+z)^{3+2\varepsilon}
  + \Omega_\Lambda(1+z)^{2\varepsilon}}$
with $\Omega_m=0.315$ fixed and $(\varepsilon,H_0)$ varied. The raw
best-fit $\varepsilon = 0.006$ is compatible with $\varepsilon=0$;
under the physical prior we therefore report this channel as a one-sided upper bound:
\begin{equation}
  \varepsilon_{\rm CC}\ :\
  \text{1\,\ensuremath{\sigma}\ upper } 0.087,\quad
  \text{2\,\ensuremath{\sigma}\ upper } 0.150.
  \label{eq:eps_cc_retained}
\end{equation}
The constraint is weak; its role is late-time consistency, not
extraction.

\paragraph{Probe 4 --- $f\sigma_8$ from RSD (one-sided consistency bound; binding above).}
Fourteen RSD measurements at $z=0.07$--$1.94$ are fit by integrating
the linear-growth ODE for the density contrast $\delta(a)$ with a
modified Poisson source driven by $G_{\rm eff}(a)$, profiling
$\sigma_8(0)$ freely. A strong $\sigma_8$--$\varepsilon$ degeneracy
leaves the raw likelihood maximum near $\varepsilon \approx -0.06$,
without a physically meaningful lower branch in the present ECT
interpretation; under the physical prior we therefore report this
channel as a one-sided upper bound:
\begin{equation}
  \varepsilon_{f\sigma_8}\ :\
  \text{1\,\ensuremath{\sigma}\ upper } 0.036,\quad
  \text{2\,\ensuremath{\sigma}\ upper } 0.120.
  \label{eq:eps_fsigma8_retained}
\end{equation}
The $1\sigma$ upper is the binding upper bound of the joint effective
band reported in \S15.5.3.

\paragraph{Probe 5 --- ISW amplitude (one-sided consistency bound under $\varepsilon\geq 0$).}
A linearised proxy $A_{\rm ISW}\approx 1+\kappa_A\,\varepsilon$, rather
than a full fit to the observed $C_\ell^{Tg}$ estimator, is compared to
Planck$+$CMB-cross-correlation data on the ISW window $z\lesssim 2$
under $\Lambda$CDM-background Limber weights. The raw likelihood
maximum lies near $\varepsilon \approx -0.007$ and remains compatible
with $\varepsilon = 0$ within a broad uncertainty; under the physical
prior we therefore report this channel as a one-sided upper bound:
\begin{equation}
  \varepsilon_{\rm ISW}\ :\
  \text{1\,\ensuremath{\sigma}\ upper } 0.043,\quad
  \text{2\,\ensuremath{\sigma}\ upper } 0.093.
  \label{eq:eps_isw_retained}
\end{equation}
Role: consistency, not extraction.

\paragraph{Effective consistency window, not a global joint likelihood.}
The retained five-probe set should not be read as a full statistically
independent joint likelihood. It is an effective multi-probe
comparison performed within a common diagnostic uniform-$\varepsilon$
layer, with partially shared background and growth ingredients across
some channels. The resulting band in \S15.5.3 is therefore an
effective consistency window under the present methodology, not a
first-principles ECT prediction and not a final global MCMC
determination. In particular, Hubble and cosmic chronometers share
background-expansion ingredients, while JWST, $f\sigma_8$, and ISW
partially overlap in their dependence on late-time growth or potential
evolution. The retained five-probe comparison should thus be read as a
structured consistency test of a common effective deformation layer,
not as a claim of strict statistical independence.

\paragraph{Bridge.}
The combined $1\sigma$ and $2\sigma$ intervals of these five effective
channels, and the resulting joint effective band under the present
methodology, are taken up in \S15.5.3.

% ==================================================================
% END §15.5.2 (rebuilt). §15.5.3--15.5.6 will be inserted here in
% subsequent steps. DO NOT REMOVE existing §15.5 below.
% ==================================================================
